\usepackage[alf,abnt-etal-cite=3,abnt-etal-list=0]{abntex2cite}
\usepackage{lmodern}	       % Usa a fonte Latin Modern
\usepackage[table]{xcolor} 
\usepackage[T1]{fontenc}		% Selecao de codigos de fonte.
\usepackage[utf8]{inputenc}		% Codificacao do documento (conversão automática dos acentos)
\usepackage{lastpage}			% Usado pela Ficha catalográfica
\usepackage{indentfirst}		% Indenta o primeiro parágrafo de cada seção.
\usepackage{color}				% Controle das cores
\usepackage{graphicx}			% Inclusão de gráficos
\usepackage{microtype} 			% para melhorias de justificação
\usepackage{graphicx}
\definecolor{color_200_200_200}{RGB}{200,200,200}

% pacotes adicionados
\usepackage{amsmath}
\usepackage{amssymb,amsfonts,amsthm}
\usepackage{setspace}
\usepackage{pdfpages} 
\usepackage{hyperref}
\usepackage{adjustbox}
% \usepackage{todonotes}

\hypersetup{
     	%pagebackref=true,
		colorlinks=true,       		% false: boxed links; true: colored links
    	linkcolor=black,          	% color of internal links
    	citecolor=black,        		% color of links to bibliography
    	filecolor=magenta,      		% color of file links
		urlcolor=blue,
		bookmarksdepth=4
}


\definecolor{blue}{RGB}{41,5,195}

% O tamanho do parágrafo é dado por:
\setlength{\parindent}{1.3cm}

% Controle do espaçamento entre um parágrafo e outro:
\setlength{\parskip}{0.2cm}  % tente também \onelineskip

% Novo list of (listings) para QUADROS

\newcommand{\quadroname}{Quadro}
\newcommand{\listofquadrosname}{Lista de quadros}

\newfloat[chapter]{quadro}{loq}{\quadroname}
\newlistof{listofquadros}{loq}{\listofquadrosname}
\newlistentry{quadro}{loq}{0}

% configurações para atender às regras da ABNT
\setfloatadjustment{quadro}{\centering}
\counterwithout{quadro}{chapter}
\renewcommand{\cftquadroname}{\quadroname\space} 
\renewcommand*{\cftquadroaftersnum}{\hfill--\hfill}

% Configuração de posicionamento padrão:
\setfloatlocations{quadro}{hbtp}
\usepackage{booktabs, multirow} % for borders and merged ranges
\usepackage{soul}% for underlines
\usepackage[table]{xcolor} % for cell colors
\usepackage{changepage,threeparttable} % for wide tables
\usepackage{array} 
\usepackage{ragged2e} 
\usepackage{makecell}
\usepackage{tabularx}

% -------------------------------------------------------------
% Controle de posicionamento de floats
% -------------------------------------------------------------
% \FloatBarrier automatico a cada \section: impede que uma figura/tabela
% declarada em uma secao "vaze" para a secao seguinte. Para barreiras
% entre subsecoes, use \FloatBarrier manualmente onde necessario.
\usepackage[section]{placeins}

% Torna o LaTeX mais tolerante ao empurrar floats: permite ate 90% da
% pagina com float no topo/rodape e apenas 5% de texto minimo (default
% 70/30/20), o que evita floats "presos" no fim de secoes longas.
\renewcommand{\topfraction}{0.90}
\renewcommand{\bottomfraction}{0.90}
\renewcommand{\textfraction}{0.05}
\renewcommand{\floatpagefraction}{0.70}

% Evita que o LaTeX estique verticalmente paginas parciais (deixadas por
% floats grandes). Melhor: alinha o texto pelo topo e mantem o espaco
% vazio no rodape sem distorcer o entrelinhamento.
\raggedbottom

% -------------------------------------------------------------
% Utilitarios de layout: quebras em identificadores longos e colunas
% -------------------------------------------------------------
% \pathtt{...} renderiza como \texttt mas permite quebra de linha em
% \_  .  /  (pontos comuns em nomes de arquivos e recursos IaC).
% Uso: \pathtt{deny\_unencrypted\_s3.rego}  <-- mesma sintaxe do \texttt
\newcommand{\pathtt}[1]{{%
  \begingroup
    \renewcommand{\_}{\char`\_\allowbreak}%
    \def\.{.\allowbreak}%
    \def\/{/\allowbreak}%
    \texttt{#1}%
  \endgroup
}}

% Tipo de coluna 'L' = paragrafo com \raggedright (evita overfull hbox
% em celulas com identificadores longos). Uso: |L{4.2cm}|L{3.1cm}|...
\newcolumntype{L}[1]{>{\raggedright\arraybackslash}p{#1}}


%\usepackage{fancyhdr}
%\pagestyle{fancy}
%\fancyhf{} % clear existing header/footer entries
% Place Page X of Y on the right-hand
% side of the footer
%\fancyfoot[R]{Página \thepage \hspace{1pt} de \pageref{LastPage}}

\usepackage{xspace}
\usepackage{tcolorbox}
\tcbuselibrary{listings, skins}

% -------------------------------------------------------------
% Trecho de codigo em linha (mantido do template original: \begin{foo}...\end{foo})
% -------------------------------------------------------------
\newtcblisting{foo}{
  listing only,
  nobeforeafter,
  after={\xspace},
  hbox,
  tcbox raise base,
  fontupper=\ttfamily,
  colback=lightgray,
  colframe=lightgray,
  size=fbox
  }

% -------------------------------------------------------------
% Blocos de codigo (listings): Terraform (HCL), YAML, Python, Dockerfile, Rego, Bash
% -------------------------------------------------------------
\usepackage{listings}

\definecolor{codebg}{RGB}{248,248,248}
\definecolor{codeframe}{RGB}{200,200,200}
\definecolor{codekw}{RGB}{0,0,180}
\definecolor{codestr}{RGB}{163,21,21}
\definecolor{codecomm}{RGB}{60,120,60}
\definecolor{codeln}{RGB}{110,110,110}

\lstdefinelanguage{HCL}{
  morekeywords={resource,module,provider,variable,output,data,locals,terraform,for,in,if,else,true,false,null,dynamic,content,depends_on,count,for_each,each,var},
  morecomment=[l]{\#},
  morecomment=[l]{//},
  morecomment=[s]{/*}{*/},
  morestring=[b]",
  sensitive=true,
}

\lstdefinelanguage{yaml}{
  morekeywords={name,on,jobs,steps,uses,with,run,env,if,needs,strategy,matrix,continue-on-error,permissions,working-directory,runs-on,apiVersion,kind,metadata,spec,containers,image,ports,resources,requests,limits,securityContext,replicas,template,selector,namespace,labels,volumes,env,valueFrom,configMap,secretKeyRef},
  morecomment=[l]{\#},
  morestring=[b]',
  morestring=[b]",
  sensitive=true,
}

\lstdefinelanguage{Rego}{
  morekeywords={package,import,default,not,some,in,contains,if,else,rule,input,data,every,as,with,true,false,null,print,sprintf},
  morecomment=[l]{\#},
  morestring=[b]",
  morestring=[b]`,
  sensitive=true,
}

\lstdefinelanguage{dockerfile}{
  morekeywords={FROM,RUN,CMD,LABEL,MAINTAINER,EXPOSE,ENV,ADD,COPY,ENTRYPOINT,VOLUME,USER,WORKDIR,ARG,ONBUILD,STOPSIGNAL,HEALTHCHECK,SHELL},
  morecomment=[l]{\#},
  morestring=[b]",
  morestring=[b]',
  sensitive=true,
}

\lstset{
  basicstyle=\ttfamily\footnotesize,
  keywordstyle=\color{codekw}\bfseries,
  stringstyle=\color{codestr},
  commentstyle=\color{codecomm}\itshape,
  numberstyle=\tiny\color{codeln},
  backgroundcolor=\color{codebg},
  frame=single,
  rulecolor=\color{codeframe},
  numbers=left,
  numbersep=6pt,
  showstringspaces=false,
  breaklines=true,
  breakatwhitespace=false,
  columns=fullflexible,
  keepspaces=true,
  tabsize=2,
  aboveskip=0.5\baselineskip,
  belowskip=0.5\baselineskip,
  literate=
    {á}{{\'a}}1 {ã}{{\~a}}1 {â}{{\^a}}1 {à}{{\`a}}1
    {é}{{\'e}}1 {ê}{{\^e}}1
    {í}{{\'i}}1
    {ó}{{\'o}}1 {õ}{{\~o}}1 {ô}{{\^o}}1
    {ú}{{\'u}}1 {ü}{{\"u}}1
    {ç}{{\c c}}1
    {Á}{{\'A}}1 {É}{{\'E}}1 {Í}{{\'I}}1 {Ó}{{\'O}}1 {Ú}{{\'U}}1
    {Ç}{{\c C}}1,
}

% Nome do quadro para blocos de codigo (padrao ABNT: "Quadro")
\renewcommand{\lstlistingname}{Quadro}
\renewcommand{\lstlistlistingname}{Lista de Quadros}